
%\needspace{4\baselineskip}
\color{uniblue}{\subsection{Контроль силового выключателя (КСВ)}}
\color{black}

\begin{enumerate}

\item
Контроль силового выключателя предназначен для определения состояния выключателя и формирования соответствующей сигнализации. 

\item
В устройстве контроль силового выключателя представлен функцией «КСВ». Функция формирует сигналы о факте наличия аварийного отключения выключателя, блокировки включения/отключения и сигнализации неисправности выключателя на основе получаемой технологической информации от силового выключателя и от других функций в составе устройства. В составе функции реализованы алгоритмы защиты ЭМВ и ЭМО.

\item
Сигнал о самопроизвольном отключении выключателя «КСВ: В самопр. отключен» формируется (при введенном в работу КСВ) при одновременном выполнении следующих условий:
\begin{itemize}
\item нет срабатывания защитных функций;
\item активное состояние выходного сигнала «КСВ: В аварийно отключен»;
\end{itemize}

\item
Сигнал о неисправности выключателя «КСВ: Неисправность В» будет формироваться (при введенном в работу КСВ) при выполнении любого из следующих условий:
\begin{itemize}
\sloppy
\item пропадание входных сигналов «ОТ цепей ЭМО1, ЭМВ» и «ОТ цепей ЭМО2» (прохождение сигналов определяется состоянием программного переключателя SGF5 «Контроль\_ОТ\_ЭМ»);
\item наличие входного сигнала «Авар.уров.изоляц. В» с подтверждением или без подтверждения от входного сигнала «Низ.уров.изоляц. В» (определяется состоянием программного переключателя SGF7 «Контроль\_сигн\_ст»);
\item наличие входного сигнала «Низ.уров.изоляц. В»;
\item наличие входного сигнала «Пружина не заведена» в течение времени, задаваемого выдержкой таймера T1 (уставка «Tгот\_в»);
\item появление входного сигнала «В: Неиспр. положение» (сигнал формируется в функциональном блоке «КА / Выключатель» по анализу положения блок-контактов выключателя);
\item активное состояние выходного сигнала «КСВ: Неиспр. цепей ЭМУ»;
\item наличие входного сигнала внешней блокировки управления выключателем «Внеш. блок. упр. В»;
\item наличие любого из входных сигналов «Работа ЭМВ», «Работа ЭМО1» или «Работа ЭМО2» в течение времени, задаваемого уставкой T3 «Тдлит\_ЭМ».
\end{itemize}

\item
Сигнал фиксации включенного положения выключателя «КСВ: РФК» формируется при наличии активного состояния сигнала «В: Включен» и отсутствии любого из условий сброса:
\begin{itemize}
    \item активное состояние сигнала «КП / УВ: Отключить»;
    \item активное состояние сигнала «Опер. отключить В»;
    \item активное состояние сигнала «Откл. от кнопки» (при введенной уставке SGF7 «Контроль\_кнопки»); 
    \item появление сигнала «Сброс» при активном состоянии сигнала «В: Отключен».
\end{itemize}

\item
Сигнал аварийного отключения выключателя «КСВ: В аварийно отключен» формируется по факту несоответствия текущего положения выключателя (входной сигнал «В: Отключен» активен) и элемента фиксации включенного положения выключателя (сигнал «КСВ: РФК» активен). Возврат элемента фиксации включенного положения выключателя (сброс сигнала «КСВ: В аварийно отключен») осуществляется при выполнении любого из следующих условий:
\begin{itemize}
\item активация входного сигнала «Сброс» при отключенном положении выключателя (входной сигнал «В: Отключен» активен);
\item активация входного сигнала «УВ: Включить» (команда оперативного включения выключателя);
\item активация входного сигнала «Откл. от кнопки» при введенном состоянии программного переключателя SGF8 «Контроль\_кнопки».
\end{itemize}

\item
Сигнал блокировки включения выключателя «КСВ: Блокировка вкл» формируется (при введенном в работу КСВ) при выполнении любого из следующих условий:
\begin{itemize}
\sloppy
\item пропадание входных сигналов «ОТ цепей ЭМО1, ЭМВ» и «ОТ цепей ЭМО2» (прохождение сигналов определяется состоянием программного переключателя SGF5 «Контроль\_ОТ\_ЭМ»);
\item появление входного сигнала «Пружина не заведена»;
\item появление входного сигнала «Авар.уров.изоляц. В» с подтверждением или без подтверждения от входного сигнала «Низ.уров.изоляц. В» (определяется состоянием программного переключателя SGF7 «Контроль\_сигн\_ст»);
\item появление входного сигнала «Низ.уров.изоляц. В» (прохождение сигнала определяется состоянием программного переключателя SGF2 «Бл\_вкл\_от\_низ.ур.КИ»);
\item активное состояние выходного сигнала «КСВ: Неиспр. цепей ЭМУ»;
\item появление входного сигнала «В: Неиспр. положение» (прохождение сигнала определяется состоянием программного переключателя SGF3 «Бл\_вкл\_от\_неисп\_пол»);
\item наличие входного сигнала «КРВ: Превышен ресурс В» (прохождение сигнала определяется состоянием программного переключателя SGF4 «Бл\_вкл\_от\_прев\_рес»);
\item наличие входного сигнала срабатывания защит;
\item наличие входного сигнала оперативного отключения выключателя «УВ: Отключить»;
\item наличие входного сигнала оперативной блокировки выключателя «ОБР В: Блокировка вкл»; %- вводится при наличии ОБР В
\item наличие любого из входных сигналов «Работа ЭМО1» или «Работа ЭМО2»;
\item наличие входного сигнала «Внеш.блок.упр.В»;
\item наличие входного сигнала «КП / УВ: Отключить»;
\item наличие входного сигнала «Опер. отключить В»;
\item наличие входного сигнала «Откл. от кнопки» при введенном состоянии программного переключателя SGF7 «Контроль\_кнопки».
\end{itemize}

\item
Сигнал блокировки отключения «КСВ: Блокировка откл» формируется (при введенном в работу КСВ) при выполнении любого из следующих условий:
\begin{itemize}
\sloppy
\item наличие входного сигнала «Авар.уров.изоляц. В» с подтверждением или без подтверждения от входного сигнала «Низ.уров.изоляц. В» (определяется состоянием программного переключателя SGF8 «Блок\_упр\_КИ\_В»);
\item наличие входного сигнала «Внеш.блок.упр.В».
\end{itemize}

\item
Сигнал неисправности цепей электромагнитов управления «КСВ: Неиспр. цепей ЭМУ» формируется при введенной функции «КСВ» в работу и состоянии программного переключателя SGF6 «Контроль\_ЭМ» отличным от «Не предусмотрено», через выдержку времени таймера T2 (уставка «Тдлит\_неисп\_ЭМУ») при выполнении любого из следующих условий:
\begin{itemize}
\item одновременном наличии или одновременном отсутствии входных сигналов «Контроль цепи ЭМВ» и «Контроль цепи ЭМО1» при состоянии программного переключателя SGF6 «Контроль\_ЭМ» равном «ЭМВ и ЭМО1» (применимо для выключателя с одним ЭМО);
\item одновременном наличии или одновременном отсутствии входных сигналов «Контроль цепи ЭМВ» и «Контроль цепи ЭМО1» при состоянии программного переключателя SGF6 «Контроль\_ЭМ» равном «ЭМВ, ЭМО1 и ЭМО2» (применимо для выключателя с двумя ЭМО);
\item одновременном наличии или одновременном отсутствии входных сигналов «Контроль цепи ЭМВ» и «Контроль цепи ЭМО2» при состоянии программного переключателя SGF6 «Контроль\_ЭМ» равном «ЭМВ, ЭМО1 и ЭМО2» (применимо для выключателя с двумя ЭМО);
\end{itemize}

\item
Сигналы срабатывания защиты электромагнитов управления («КСВ: Защита ЭМВ», «КСВ: Защита ЭМО1», «КСВ: Защита ЭМО2») формируются при введенной функции «КСВ» и наличии сигналов «Работа ЭМВ», «Работа ЭМО1» и «Работа ЭМО2» в течение времени, задаваемого уставкой T3 «Тдлит. ЭМ» для каждой цепи отдельно.

\item
Ввод функции в работу на этапе параметрирования устройства осуществляется программным переключателем SGF1 «Ввод\_функции». Информация о введенной функции отображается выходным сигналом «КСВ: Ввод».
\item
Функция может быть выведена из работы оперативно путем активации сигнала «ОВ КСВ», что характеризуется активным состоянием сигнала «КСВ: Оперативный вывод» на выходе функции.

\item
Алгоритм работы функции выполнен в соответствии с рисунком \ref{ksv:img1}. В таблице \ref{ksv:tbl1} приведены параметры, необходимые для настройки функции.

\ifboolexpr{bool {isVV}}
{
\vspace{3mm}
\begin{figure}[H]
\centering
\includegraphics[width=0.9\textwidth,height=0.9\textheight,keepaspectratio]{\fbpath/314. КСВ ВВ/_latex/img/ksv.pdf}
\captionof{figure}{Функциональная схема логики КСВ}
\label{ksv:img1}
\end{figure}
}
{
\vspace{3mm}
\begin{figure}[H]
\centering
\includegraphics[width=0.9\textwidth,height=0.9\textheight,keepaspectratio]{\fbpath/314. КСВ ВВ/_latex/img/svksv.pdf}
\captionof{figure}{Функциональная схема логики КСВ}
\label{ksv:img1}
\end{figure}
}

\end{enumerate}
